\documentclass[12pt]{article}
\usepackage[utf8]{inputenc}
\usepackage[T1]{fontenc}
\usepackage{hyperref}
\title{Tangut Language Resources for K-12 Curriculum Development}
\author{Compiled by Claude Sonnet 5 (claude-sonnet-5) for Andr\'as Kornai}
\date{}
\begin{document}
\maketitle

\section{Introduction}
Tangut (Xixia) is an extinct Sino-Tibetan language of the Western Xia state (11th to 13th centuries), written in its own logographic script. It has no native speakers, and no K-12 programme was found. Resources are scholarly dictionaries and corpora; the script is encoded in Unicode (6,125 characters in Unicode 9.0, 2016), which makes digital materials possible.

\section{Primary Educational Resources}
No school primers were found. The Concise Online English-Tangut Dictionary (LMU Munich, World Philology) is the most accessible entry point, giving characters with Chinese and English meanings and transcriptions: \url{https://www.worldphilology.gwi.uni-muenchen.de/?p=3780}

\section{Linguistic and Reference Materials}
\begin{itemize}
\item Corpus Textuum Tangutorum series (BabelStone): \url{https://babelstone.co.uk/CTT/index.html}
\item Gerard Clauson's Skeleton Tangut (Hsi Hsia) Dictionary, facsimile edition (Evertype): \url{https://evertype.com/books/ctt1-clauson.html}
\item Li Fanwen's Tangut-Chinese dictionary (background): \url{https://www.wikipedia.org/wiki/Li_Fanwen}
\item Unicode proposal for the Tangut script (PDF): \url{https://www.unicode.org/L2/L2007/07143-tangut.pdf}
\item Omniglot on the Tangut script: \url{https://www.omniglot.com/writing/tangut.htm}
\item Kychanov and Arakawa, Tangut dictionary (Kyoto, 2006) and Han Xiaomang's nine-volume dictionary (2021), described at: \url{https://teahouse.buddhistdoor.net/?p=17774}
\end{itemize}

\section{Religious and Community Resources}
Most surviving texts are Buddhist translations. No modern community organisation was found in the search results.

\section{Digital and Audio Resources}
Unicode encoding and the online dictionary above; no audio exists, since the pronunciation is only reconstructed.

\section{Curriculum Development Context}
A K-12 unit could be an enrichment module on writing systems and Silk Road history, using the script and a small vocabulary from the online dictionary. Full language instruction is not realistic. Availability and licensing of the print works was not checked.

\end{document}